% !PRINTED 591
\weberpage{591}
\begin{center}
\S.\,161.\\[4pt]
\gesperrt{Primitive Roots of the Prime Ideals.}
\end{center}

From the formulae (4), (5) of the preceding section it follows, having regard to \S.\,150, that for every integer $\eta$ of the field $R$
\[
\eta^P \equiv \eta \pmod{\mathfrak{P}},
\]
hence also
\weq{1}{\eta^P \equiv \eta \pmod{\mathfrak{p}},}
and for every number $\omega$ in $\mathfrak{o}$
\weq{2}{\omega^{P^f} \equiv \omega \pmod{\mathfrak{p}}.}
The number of integers in $R$ incongruent with respect to $\mathfrak{P}$ is $P$, and the number of numbers in $\Omega$ incongruent with respect to $\mathfrak{p}$ is $P^f$.

We have already shown earlier (\S.\,150) that for every prime ideal $\mathfrak{p}$ there are \gesperrt{primitive roots} $\gamma$, that is, numbers in $\Omega$ which satisfy the congruence
\weq{3}{\gamma^{P^f-1} \equiv 1 \pmod{\mathfrak{p}}}
and at the same time have the property that no lower power with positive exponent becomes congruent to unity. Then every integer in $\Omega$ not divisible by $\mathfrak{p}$ is congruent to one and only one of the powers of $\gamma$
\weq{4}{1, \gamma, \gamma^2, \ldots, \gamma^{P^f-2}}
with respect to the modulus $\mathfrak{p}$.

Every integer $\Theta$ in $\Omega$ satisfies an equation of degree at most $n$, whose coefficients are integers in $R$.

This equation is at the same time a congruence with respect to the modulus $\mathfrak{p}$, and among all such congruences there will be one of \gesperrt{lowest possible degree}
\[
F(\Theta) \equiv 0 \pmod{\mathfrak{p}}
\]
in which, moreover, we may assume the coefficient of the highest power of $\Theta$ to be equal to 1. For if the highest coefficient (indivisible by $\mathfrak{p}$, and consequently by $\mathfrak{P}$) is $\eta_0$, then we can always satisfy the congruence
\[
\eta_0 \eta \equiv 1 \pmod{\mathfrak{P}}
\]
by an integer $\eta$ in $R$, and have then only to multiply the function $F$ by this $\eta$. We can
%CONT>>
